% =============================================================================
% Canonical template for 0-lifecycle/0-seed/0-seed.tex
% Standalone-compilable. The skill owns this body; the preamble matches the
% scaffold (article + geometry + booktabs + hyperref + title/date).
% Section order is fixed: Seed Question -> Motivations -> Tentative Claim Shape
% -> Current Evidence Status -> Open Evidence Needs -> Promotion Gate -> Kill
% Criteria. Layout: paragraph banner + one sentence per line + %% ---- Pn.Sm ----
% tags (paragraph IDs are file-local, restart at P1).
% =============================================================================
\documentclass[11pt]{article}
\usepackage[margin=1in]{geometry}
\usepackage{booktabs}
\usepackage{hyperref}
\title{Seed: <Paper Title>}
\date{<YYYY-MM-DD>}
\begin{document}
\maketitle

\section*{Seed Question}

% =========================================================
% Para [seed.question] Setup -- the paper-shaped question
% =========================================================
%% ---- P1.S1 ----
<The single paper-shaped question this seed exists to answer.>

\section*{Motivations}

% =========================================================
% Para [seed.motivation] Rationale -- why interesting, and to whom
% =========================================================
%% ---- P2.S1 ----
<Why is this interesting: the puzzle, the gap, or the surprise.>
%
%% ---- P2.S2 ----
<What makes the angle novel or feasible now.>
%
%% ---- P2.S3 ----
<To whom is it interesting: name the audiences (clinicians, payers/policymakers, a research community) and why each cares.>

\section*{Tentative Claim Shape}

% =========================================================
% Para [seed.shape] Rationale -- what the paper may argue, as hypotheses
% =========================================================
%% ---- P3.S1 ----
<The core hypothesis, phrased as a hypothesis, not a finding.>
%
%% ---- P3.S2 ----
<Any secondary shapes the paper may take, still hedged.>

\section*{Current Evidence Status}

% =========================================================
% Para [seed.evidence] Status -- the evidence that already exists
% =========================================================
%% ---- P4.S1 ----
<Task/probe/discovery/insight state. Say plainly when this is not yet a paper seed.>

\section*{Open Evidence Needs}

% =========================================================
% Para [seed.needs] Action -- what to get next, each with a route
% =========================================================
%% ---- P5.S1 ----
<Need, with a route: probe / discover / task / insight.>

\section*{Promotion Gate}

% =========================================================
% Para [seed.gate] Decision -- concrete conditions to advance to pitch
% =========================================================
%% ---- P6.S1 ----
<What must be true to promote this to an active paper seed, then to pitch.>

\section*{Kill Criteria}

% =========================================================
% Para [seed.kill] Decision -- what would make this paper not worth pursuing
% =========================================================
%% ---- P7.S1 ----
<What evidence would make this paper not worth pursuing.>

\end{document}
